% v2 -- rewritten
% Seções:
%   5.1 - Modelagem 3D do alinhador ortodôntico
%   5.2 - Aplicação de cargas e simulação FEM
%   5.3 - Validação dos resultados
%   5.4 - Projeto prático: simular 3mm de movimentação
\destaque{Nivel-alvo N1 — modifique scripts de simulacao FEM com Open3D
para analisar tensoes ortodônticas no ambito do Framework SUS-Twin.}

\label{ch:ortodontia}

\section{Modelagem 3D do alinhador ortodôntico}

O alinhador\footnote{Aparelho removivel transparente de polimero termoplastico
que movimenta os dentes por meio de forcas controladas, substituto moderno do
aparelho fixo metalico.} ortodôntico e representado por uma malha poligonal\footnote{Conjunto de poligonos (triangulos, na pratica) que formam a superficie
externa de um objeto 3D, similar a uma ``pele digital'' do modelo fisico.}
gerada a partir de dados de escaneamento intraoral (IOS). Cada malha contem
vertices\footnote{Pontos no espaco 3D (coordenadas x,y,z) que definem os cantos
dos triangulos.} e faces\footnote{Triangulos formados pela conexao de tres
vertices; o conjunto de todas as faces compoe a superficie visivel do modelo.}
que descrevem a geometria do dente e do alinhador. O codigo a seguir carrega
e visualiza um arquivo STL com a biblioteca Open3D:

\begin{lstlisting}[language=Python,caption={Carregamento e visualizacao de
malha 3D de alinhador ortodôntico com Open3D}]
import open3d as o3d
import numpy as np

# Carrega a malha do alinhador a partir de arquivo STL
malha = o3d.io.read_triangle_mesh("alinhador.stl")
malha.compute_vertex_normals()

# Extrai vertices e faces para inspecao
vertices = np.asarray(malha.vertices)
faces = np.asarray(malha.triangles)
print(f"Vertices: {len(vertices)}, Faces: {len(faces)}")

# Visualizacao interativa
o3d.visualization.draw_geometries([malha],
    window_name="Alinhador - SUS-Twin",
    width=800, height=600)
\end{lstlisting}

A matriz de vertices alimenta diretamente o solver de elementos finitos no
Framework SUS-Twin, que opera sobre a mesma representacao geometrica para
calcular deformacoes sob carga.

\section{Aplicacao de cargas e simulacao FEM}

O Metodo dos Elementos Finitos\footnote{Tecnica numerica que divide um objeto
complexo em milhares de pequenas partes (elementos) para calcular, uma a uma,
como ele se deforma sob forcas aplicadas.} (FEM) discretiza a malha em
elementos tetraedricos e resolve equacoes de equilibrio para cada no.
O modulo de Young\footnote{Medida de rigidez do material: quanto maior o valor,
mais o material resiste a deformacao elastica. Acrilico dental tem ~2 GPa.}
do polimero (PGA, ~2 GPa) e a condicao de contorno\footnote{Restricao imposta
a uma regiao do modelo — neste caso, fixa-se a borda gengival do alinhador
para simular seu encaixe real nos dentes.} (fixacao na margem gengival)
definem o comportamento mecanico. Uma forca de 0,5~N\footnote{Equivalente ao
peso de uma moeda de 50 centavos (7,0~g multiplicado pela gravidade);
considerada forca ortodontica leve e biologica.} e aplicada na face vestibular:

\begin{lstlisting}[language=Python,caption={Pipeline FEM conceitual para
alinhador ortodôntico}]
# Propriedades do material (PGA)
E = 2.0e9         # Modulo de Young [Pa]
nu = 0.35         # Coeficiente de Poisson

# Condicao de contorno: nos da margem gengival fixos
bc_fixos = identificar_nos_margem(malha)
for no in bc_fixos:
    aplicar_deslocamento_zero(no)

# Forca vestibular de 0.5 N
forca = np.array([0.0, 0.5, 0.0])  # direcao vestibulo-lingual
nos_carga = identificar_nos_vestibulares(malha, regiao="incisivo")
for no in nos_carga:
    aplicar_forca(no, forca / len(nos_carga))

# Chama o solver FEM (calcula deslocamentos em cada no)
deslocamentos = solver_fem(malha, E, nu, bc_fixos, nos_carga)
print(f"Deslocamento maximo: {np.max(np.linalg.norm(
    deslocamentos, axis=1)):.4f} mm")
\end{lstlisting}

O Framework SUS-Twin abstrai este pipeline em uma unica chamada de API,
permitindo iterar rapidamente entre geometria, material e condicoes de carga
sem reescrever o solver.

\section{Validacao dos resultados}

A analise dos resultados compara o deslocamento simulado com valores
esperados na literatura clinica (0,2--0,5~mm para forcas leves). A tensao de
von Mises\footnote{Indicador unico que combina as tres direcoes de tensao em
um numero escalar; quando ultrapassa o limite de escoamento do material, o
objeto deforma-se permanentemente.} maxima identifica regioes de risco de
falha estrutural do alinhador:

\begin{lstlisting}[language=Python,caption={Calculo de deslocamento e tensao
de von Mises para validacao}]
from utils_fem import calcular_von_mises

# Calcula magnitude do deslocamento por no
mag = np.linalg.norm(deslocamentos, axis=1)
idx_max = np.argmax(mag)
print(f"Desloc. max.: {mag[idx_max]:.3f} mm "
      f"no vertice {idx_max}")

# Tensao de von Mises nos elementos
tensoes = calcular_von_mises(malha, deslocamentos, E, nu)
idx_stress = np.argmax(tensoes)
print(f"Tensao max.: {tensoes[idx_stress]:.2f} MPa")
\end{lstlisting}

O MPa\footnote{Megapascal: unidade de pressao;
1~MPa = 10~kgf/cm².} e a unidade padrao para tensao mecanica em engenharia.

Se o deslocamento calculado ficar dentro da faixa esperada e a tensao de
von Mises nao exceder o limite de escoamento do PGA (~50~MPa), a simulacao
e considerada validada. O Framework SUS-Twin inclui metricas automaticas de
conformidade biomecanica para cada cenario simulado.

\section{Projeto pratico: simular 3~mm de movimentacao}

O projeto guiado a seguir aplica todo o pipeline do Framework SUS-Twin para
simular 3~mm de translacao vestibular de um incisivo central:

\begin{itemize}
    \item \textbf{Obter o STL\footnote{Formato de arquivo 3D que armazena a
    superficie do objeto como uma lista de triangulos; padrao industrial para
    impressoras 3D e softwares CAD.}:} escaneie o modelo dental com scanner
    intraoral ou baixe um arquivo STL de repositorio aberto.
    \item \textbf{Definir o material:} configure modulo de Young $E = 2,0$
    GPa, coeficiente de Poisson $\nu = 0,35$ e espessura do alinhador
    0,6~mm (valor tipico de mercado).
    \item \textbf{Executar a simulacao:} aplique uma sequencia de 5 cargas
    incrementais de 0,1~N a 0,5~N, registrando o deslocamento a cada passo.
    \item \textbf{Validar o deslocamento:} confira se a resultante e de
    $3,0 \pm 0,3$~mm e se a tensao de von Mises maxima nao ultrapassa
    45~MPa.
\end{itemize}

Ao final, o aluno tera executado o ciclo completo de simulacao biomecanica
ortodontica — da malha bruta a validacao clinica — exatamente como preconizado
pelo Framework SUS-Twin para a criacao de gemeos digitais odontologicos
reprodutiveis e auditaveis.
